\resume
\selectlanguage{french}
\begin{abstract}
Ce mémoire de fin de cycle présente une solution numérique conçue pour répondre aux insuffisances constatées dans la gestion des travaux dirigés destinés aux candidats des examens du CEP, du BEPC et du Baccalauréat au sein de la commune de Cotonou. Dans la grande majorité des services municipaux concernés, le suivi de ces activités repose encore sur des supports papier ou des tableurs non structurés, engendrant des risques d'erreurs, des pertes d'informations, une traçabilité insuffisante et une prise de décision ralentie.

Pour répondre à cette problématique, nous avons conçu une application web centralisée et sécurisée dédiée à la modernisation de ce processus administratif. La plateforme permet la gestion structurée des sessions de travaux dirigés, le suivi rigoureux des données financières et des paiements --- qu'ils soient bancaires ou électroniques --- ainsi que la génération d'états récapitulatifs et de rapports détaillés à des fins décisionnelles. Un tableau de bord interactif offre une visibilité en temps réel sur l'ensemble des opérations. L'application intègre également un système d'authentification robuste et une gestion granulaire des droits d'accès, garantissant la confidentialité et l'intégrité des données traitées.

La solution est accessible depuis tout appareil connecté --- ordinateur, tablette ou smartphone --- sans nécessiter d'installation préalable, ce qui la rend parfaitement adaptée aux réalités des administrations locales béninoises.

Sur le plan méthodologique, la conception du système a été conduite en langage UML à travers les diagrammes de cas d'utilisation, de classes et de séquence. Du point de vue technologique, le développement frontend s'appuie sur HTML, CSS et JavaScript, tandis que le backend est réalisé avec le framework Laravel associé à une base de données MySQL, assurant robustesse, sécurité et maintenabilité de l'ensemble de l'application.

\paragraph{}
\textbf{Mots clés} : travaux dirigés, gestion administrative, paiements, digitalisation, commune de Cotonou, examens scolaires, application web.
\end{abstract}

\newpage
\thispagestyle{empty}
\selectlanguage{english}
\addcontentsline{toc}{chapter}{Abstract}
\begin{abstract}
This end-of-cycle thesis presents a digital solution designed to address the shortcomings observed in the management of tutorial sessions (travaux dirigés) for candidates of the CEP, BEPC, and Baccalaureate examinations within the municipality of Cotonou. In the vast majority of relevant municipal services, the tracking of these activities still relies on paper-based systems or unstructured spreadsheets, leading to risks of errors, loss of information, insufficient traceability, and delayed decision-making.

To address this issue, we have designed a centralized and secure web application dedicated to the modernization of this administrative process. The platform enables structured management of tutorial sessions, rigorous tracking of financial data and payments --- whether via bank transfer or electronic means --- as well as the generation of summary statements and detailed reports for decision-making purposes. An interactive dashboard provides real-time visibility over all operations. The application also integrates a robust authentication system and granular access control management, ensuring the confidentiality and integrity of the processed data.

The solution is accessible from any connected device --- computer, tablet, or smartphone --- without requiring prior installation, making it perfectly suited to the realities of local administrations in Benin.

Methodologically, the system design was conducted using the UML language through use case, class, and sequence diagrams. Technologically, the frontend development is based on HTML, CSS, and JavaScript, while the backend is implemented with the Laravel framework coupled with a MySQL database, ensuring the robustness, security, and maintainability of the entire application.

\paragraph{}
\textbf{Keywords} : tutorial sessions, administrative management, payments, digitalization, municipality of Cotonou, school examinations, web application.
\end{abstract}